\let\negmedspace\undefined
\let\negthickspace\undefined
\documentclass[journal,12pt,onecolumn]{IEEEtran}
\usepackage{cite}
\usepackage{amsmath,amssymb,amsfonts,amsthm}
\usepackage{algorithmic}
\usepackage{graphicx}
\graphicspath{{./figs/}}
\usepackage{textcomp}
\usepackage{xcolor}
\usepackage{txfonts}
\usepackage{listings}
\usepackage{enumitem}
\usepackage{mathtools}
\usepackage{gensymb}
\usepackage{comment}
\usepackage{caption}
\usepackage[breaklinks=true]{hyperref}
\usepackage{tkz-euclide} 
\usepackage{listings}
\usepackage{gvv}                                        
%\def\inputGnumericTable{}                                 
\usepackage[latin1]{inputenc}     
\usepackage{xparse}
\usepackage{color}                                            
\usepackage{array}
\usepackage{longtable}                                       
\usepackage{calc}                                             
\usepackage{multirow}
\usepackage{multicol}
\usepackage{hhline}                                           
\usepackage{ifthen}                                           
\usepackage{lscape}
\usepackage{tabularx}
\usepackage{array}
\usepackage{float}
\newtheorem{theorem}{Theorem}[section]
\newtheorem{problem}{Problem}
\newtheorem{proposition}{Proposition}[section]
\newtheorem{lemma}{Lemma}[section]
\newtheorem{corollary}[theorem]{Corollary}
\newtheorem{example}{Example}[section]
\newtheorem{definition}[problem]{Definition}
\newcommand{\BEQA}{\begin{eqnarray}}
\newcommand{\EEQA}{\end{eqnarray}}
\newcommand{\define}{\stackrel{\triangle}{=}}
\theoremstyle{remark}
\newtheorem{rem}{Remark}

\begin{document}

\title{Eigen Values and Vectors}
\author{ee25btech11056 - Suraj.N}
\maketitle
\renewcommand{\thefigure}{\theenumi}
\renewcommand{\thetable}{\theenumi}

\begin{document}

\textbf{Question :} What can we conclude if a matrix $\vec{A}$ is symmetric?\\

\textbf{Proof :}\\

\textbf{Case 1 : $\vec{A}$ is real and symmetric}\\

Let $\vec{A} \in \mathbb{R}^{n \times n}$ be a real symmetric matrix, i.e.,
\begin{align}
\vec{A} = \vec{A}^\top
\end{align}

Suppose $\lambda$ is an eigenvalue of $\vec{A}$ with eigenvector $\vec{v} \ne \vec{0}$:
\begin{align}
\vec{A}\vec{v} = \lambda \vec{v}
\end{align}

Taking the conjugate transpose of $\vec{v}$ and multiplying on the left:
\begin{align}
\vec{v}^* \vec{A}\vec{v} = \lambda \vec{v}^* \vec{v}
\end{align}

Since $\vec{A}$ is real symmetric, $\vec{A} = \vec{A}^* = \vec{A}^\top$, therefore:
\begin{align}
\vec{v}^* \vec{A}\vec{v} = (\vec{A}\vec{v})^* \vec{v} = (\lambda \vec{v})^* \vec{v} = \bar{\lambda} \, \vec{v}^* \vec{v}
\end{align}

Equating both expressions for $\vec{v}^* \vec{A}\vec{v}$:
\begin{align}
\lambda \vec{v}^* \vec{v} = \bar{\lambda} \, \vec{v}^* \vec{v}\\
(\lambda- \bar{\lambda})\vec{v}^* \vec{v} = \vec{0}
\end{align}

Since $\vec{v} \ne \vec{0}$, $\vec{v}^*\vec{v} > 0$. Hence:
\begin{align}
\lambda = \bar{\lambda}
\end{align}

Therefore, every eigenvalue $\lambda$ of $\vec{A}$ is real.\\

\textbf{Eigenvectors corresponding to distinct eigenvalues are orthogonal}\\

Let

\begin{align}
\vec{A}\vec{u} = \alpha \vec{u}\\
\vec{A}\vec{w} = \beta \vec{w}
\end{align}

Where $\alpha \ne \beta$ are distinct eigenvalues.Multiplying the first equation on the left by $\vec{w}^\top$:

\begin{align}
\vec{w}^\top \vec{A} \vec{u} = \alpha \vec{w}^\top \vec{u}
\end{align}

Since $\vec{A}$ is symmetric, $\vec{A} = \vec{A}^\top$. Hence:
\begin{align}
\vec{w}^\top \vec{A}\vec{u} = (\vec{A}\vec{w})^\top \vec{u} = (\beta \vec{w})^\top \vec{u} = \beta \vec{w}^\top \vec{u}
\end{align}

\pagebreak

Equating both forms:
\begin{align}
\alpha \vec{w}^\top \vec{u} = \beta \vec{w}^\top \vec{u}
\end{align}

Simplifying:
\begin{align}
(\alpha - \beta) \vec{w}^\top \vec{u} = 0
\end{align}

Since $\alpha \ne \beta$, it follows that:
\begin{align}
\vec{w}^\top \vec{u} = 0
\end{align}

Hence, eigenvectors corresponding to distinct eigenvalues are orthogonal.\\

\textbf{Proving that Eigen Vector is real}

Suppose $\lambda \in \mathbb{R}$ is an eigenvalue of $\vec{A}$ with a (possibly complex) eigenvector $\vec{v} \in \mathbb{C}^n$, satisfying
\begin{align}
\vec{A}\vec{v} = \lambda \vec{v}
\end{align}

\textbf{Considering the conjugate eigenvector}

Since $\vec{A}$ is real, taking the complex conjugate of both sides gives
\begin{align}
\vec{A}\bar{\vec{v}} = \bar{\lambda} \bar{\vec{v}}
\end{align}

But for a real symmetric matrix, $\lambda = \bar{\lambda}$ (as proven earlier). Hence,
\begin{align}
\vec{A}\bar{\vec{v}} = \lambda \bar{\vec{v}}
\end{align}

Thus, $\bar{\vec{v}}$ is also an eigenvector corresponding to the same eigenvalue $\lambda$.

\bigskip

\textbf{Constructing real vectors from $\vec{v}$ and $\bar{\vec{v}}$}

Define
\begin{align}
\vec{u} &= \vec{v} + \bar{\vec{v}}\\
\vec{w} &= i(\vec{v} - \bar{\vec{v}})
\end{align}

Both $\vec{u}$ and $\vec{w}$ are \textbf{real vectors}, because $\vec{u}$ is the sum of a vector and its conjugate, and $\vec{w}$ is $i$ times their difference.

\bigskip

\textbf{Applying $\vec{A}$ to $\vec{u}$ and $\vec{w}$}

Since $\vec{A}$ is real and linear,
\begin{align}
\vec{A}\vec{u} &= \vec{A}(\vec{v} + \bar{\vec{v}}) = \vec{A}\vec{v} + \vec{A}\bar{\vec{v}} = \lambda \vec{v} + \lambda \bar{\vec{v}} = \lambda (\vec{v} + \bar{\vec{v}}) = \lambda \vec{u}
\end{align}

Similarly,
\begin{align}
\vec{A}\vec{w} &= \vec{A}\big(i(\vec{v} - \bar{\vec{v}})\big) = i(\vec{A}\vec{v} - \vec{A}\bar{\vec{v}}) = i(\lambda \vec{v} - \lambda \bar{\vec{v}}) = \lambda i(\vec{v} - \bar{\vec{v}}) = \lambda \vec{w}
\end{align}

Hence we get 

\begin{align}
  \vec{A}\vec{u}=\lambda\vec{u}\\
  \vec{A}\vec{w}=\lambda\vec{w}
\end{align}
\bigskip

\textbf{Real eigenvectors}

If $\vec{u} \ne \vec{0}$, then $\vec{u}$ is a real eigenvector corresponding to eigenvalue $\lambda$.  
If $\vec{w} \ne \vec{0}$, then $\vec{w}$ is a real eigenvector corresponding to eigenvalue $\lambda$.  

\textbf{Conclusion :}
For any real symmetric matrix $\vec{A}$,
\begin{enumerate}
  \item All eigenvalues of $\vec{A}$ are \textbf{real}.
  \item Eigenvectors corresponding to distinct eigenvalues are \textbf{orthogonal}.
  \item For every eigenvalue $\lambda$ of a real symmetric matrix $\vec{A}$, there exists a \textbf{real eigenvector}.  
Hence, the eigenvectors of real symmetric matrices are real. 
\end{enumerate}

\pagebreak 

\textbf{Case 2 : $\vec{A}$ is complex and symmetric}\\

Let $\vec{A} \in \mathbb{C}^{n \times n}$ be a complex symmetric matrix satisfying
\begin{align}
\vec{A} = \vec{A}^\top
\end{align}

Note that $\vec{A}$ is not necessarily Hermitian since in general $\vec{A} \ne \vec{A}^*$.

\begin{align}
\vec{A} \ne \vec{A}^*
\end{align}

\textbf{Eigenvalues of $\vec{A}$ may be complex}

Let $\lambda \in \mathbb{C}$ be an eigenvalue of $\vec{A}$ with eigenvector $\vec{v} \ne \vec{0}$ such that
\begin{align}
\vec{A}\vec{v} = \lambda \vec{v}
\end{align}

Define scalar z as :

\begin{align}
z = \vec{v}^* \vec{A} \vec{v}  
\end{align}

Taking complex conjugate transpose on both sides ,

\begin{align}
z^* = \vec{v}^* \vec{A}^* \vec{v}   
\end{align}

If $\vec{A} = \vec{A}^*$ , then 

\begin{align}
  z = z^* \\
  z = \bar{z} \text{   as z is a scalar}
\end{align}

Therefore z is a real scalar value \\

For Hermitian matrices ($\vec{A} = \vec{A}^*$), the quantity $\vec{v}^* \vec{A} \vec{v}$ is always real.  
However, for complex symmetric matrices, $\vec{A} \ne \vec{A}^*$, so

\begin{align}
\vec{v}^* \vec{A} \vec{v} \text{ may not be real.}
\end{align}

The Rayleigh quotient is
\begin{align}
\frac{\vec{v}^* \vec{A} \vec{v}}{\vec{v}^* \vec{v}} = \lambda
\end{align}

Since $\vec{v}^* \vec{A} \vec{v}$ may be complex, $\lambda$ may also be complex.

\textbf{Eigenvectors corresponding to distinct eigenvalues need not be orthogonal}

For complex symmetric matrices, we have only $\vec{A} = \vec{A}^\top$, so the Hermitian symme
try property does not hold.

Thus, for two eigenpairs $(\lambda_1, \vec{v}_1)$ and $(\lambda_2, \vec{v}_2)$ with $\lambda_1 \ne \lambda_2$, we \textbf{cannot guarantee}
\begin{align}
\vec{v}_1^* \vec{v}_2 = 0
\end{align}

Hence, eigenvectors corresponding to distinct eigenvalues \textbf{need not be orthogonal and are not guaranteed to be real}.

\textbf{Conclusion :}  
For a complex symmetric matrix $\vec{A}$:
\begin{enumerate}
  \item The eigenvalues \textbf{can be complex}.
  \item The eigenvectors corresponding to distinct eigenvalues are \textbf{not necessarily orthogonal and need not be real}.
\end{enumerate}
 
But if $\vec{A} = \vec{A}^*$ then we get the same result as \textbf{case 1}

\pagebreak

\textbf{Examples to Support the Proof}

\textbf{Example 1 : Real Symmetric Matrix}

(\textbf{a}) Let
\begin{align}
\vec{A} = \myvec{2 & -1 & 0\\ -1 & 2 & -1\\ 0 & -1 & 2}
\end{align}

Clearly, $\vec{A} = \vec{A}^\top$, hence $\vec{A}$ is real symmetric.

\textbf{Finding Eigenvalues}

The characteristic equation is:
\begin{align}
\mydet{\vec{A} - \lambda \vec{I}} = 0
\end{align}

Substituting $\vec{A}$,
\begin{align}
\mydet{
\myvec{
2-\lambda & -1 & 0\\
-1 & 2-\lambda & -1\\
0 & -1 & 2-\lambda
}
} = 0
\end{align}

\begin{align}
\myvec{
2-\lambda & -1 & 0\\
-1 & 2-\lambda & -1\\
0 & -1 & 2-\lambda
}
\xleftrightarrow{R_2 \to R_2 + \frac{1}{2-\lambda}R_1}
\myvec{
2-\lambda & -1 & 0\\
0 & \frac{(2-\lambda)^2 - 1}{2-\lambda} & -1\\
0 & -1 & 2-\lambda
}
\xleftrightarrow{R_3 \to R_3 + \frac{(2-\lambda)}{(2-\lambda)^2 - 1}R_2}
\myvec{
2-\lambda & -1 & 0\\
0 & \frac{(2-\lambda)^2 - 1}{2-\lambda} & -1\\
0 & 0 & \frac{(2-\lambda)^3 - 4(2-\lambda)}{(2-\lambda)^2 - 1}
}
\end{align}

\begin{align}
(2-\lambda)\Big(\frac{(2-\lambda)^2 - 1}{2-\lambda}\Big)\Big(\frac{(2-\lambda)^3 - 4(2-\lambda)}{(2-\lambda)^2 - 1}\Big) = 0
\end{align}

Simplifying,
\begin{align}
(2-\lambda)\big((2-\lambda)^2 - 2\big) = 0
\end{align}

Hence, the eigenvalues are:
\begin{align}
\lambda_1 = 2,\ \lambda_2 = 2 + \sqrt{2},\ \lambda_3 = 2 - \sqrt{2}
\end{align}

All eigenvalues are real, confirming positive definiteness.\\

\textbf{Eigenvectors:}
Solving $(\vec{A}-\lambda_i\vec{I})\vec{v}_i=0$ gives:
\begin{align}
\vec{v}_1 = \myvec{1\\0\\-1},\ 
\vec{v}_2 = \myvec{1\\\sqrt{2}\\1},\ 
\vec{v}_3 = \myvec{1\\-\sqrt{2}\\1}
\end{align}

Thus, eigenvectors are real and mutually orthogonal.

\pagebreak

\textbf{Example 2 : Complex Symmetric Matrix}

(\textbf{a}) Let
\begin{align}
\vec{A} = \myvec{1 & 1+i\\ 1+i & 1}
\end{align}

Here, $\vec{A} = \vec{A}^\top$ but $\vec{A} \ne \vec{A}^*$, so $\vec{A}$ is complex symmetric but not Hermitian.

\textbf{Finding Eigenvalues}

\begin{align}
\mydet{\vec{A} - \lambda \vec{I}} = 0
\end{align}

\begin{align}
\mydet{
\myvec{
1-\lambda & 1+i\\
1+i & 1-\lambda
}
} = 0
\end{align}

\begin{align}
\mydet{
1-\lambda & 1+i\\
1+i & 1-\lambda
}
\xleftrightarrow{R_2 \to R_2 - \frac{1+i}{1-\lambda} R_1}
\mydet{
1-\lambda & 1+i\\
0 & (1-\lambda) - \frac{(1+i)^2}{1-\lambda}
}
\end{align}

\begin{align}
(1-\lambda)^2 - (1+i)^2 = 0
\end{align}

\begin{align}
(1-\lambda)^2 - 2i = 0
\end{align}

\begin{align}
1-\lambda = \pm \sqrt{2i} \quad \Rightarrow \quad \lambda = 1 \mp \sqrt{2i}
\end{align}

Hence eigenvalues are:
\begin{align}
\lambda_1 = 1 - \sqrt{2i}, \quad \lambda_2 = 1 + \sqrt{2i}
\end{align}

\textbf{Eigenvectors:}

\textbf{Eigenvalue } \(\lambda_1 = 1 - \sqrt{2i}\):

\begin{align}
\vec{v}_1 = \myvec{1\\ -\,\frac{\sqrt{2i}}{1+i}}
\end{align}

\textbf{Eigenvalue } \(\lambda_2 = 1 + \sqrt{2i}\):

\begin{align}
\vec{v}_2 = \myvec{1\\ \frac{\sqrt{2i}}{1+i}}
\end{align}

\textbf{Summary:}
\begin{align}
\lambda_1 &= 1 - \sqrt{2i}, \quad \vec{v}_1 = \myvec{1\\ -\,\frac{\sqrt{2i}}{1+i}}\\
\lambda_2 &= 1 + \sqrt{2i}, \quad \vec{v}_2 = \myvec{1\\ \frac{\sqrt{2i}}{1+i}}
\end{align}

The Eigen vectors are \textbf{complex and not orthogonal}

\pagebreak

(\textbf{b}) Let
\begin{align}
\vec{A} = \myvec{1 & i & 0\\ i & 1 & i\\ 0 & i & 1}
\end{align}

Here, $\vec{A} = \vec{A}^\top$ but $\vec{A} \ne \vec{A}^*$, so $\vec{A}$ is complex symmetric but not Hermitian.

\textbf{Finding Eigenvalues}

\begin{align}
\mydet{\vec{A} - \lambda \vec{I}} = 0
\end{align}

\begin{align}
\mydet{
\myvec{
1-\lambda & i & 0\\
i & 1-\lambda & i\\
0 & i & 1-\lambda
}
} = 0
\end{align}

\begin{align}
\myvec{
1-\lambda & i & 0\\
i & 1-\lambda & i\\
0 & i & 1-\lambda
}
\xleftrightarrow{R_2 \to R_2 - \frac{i}{1-\lambda} R_1}
\myvec{
1-\lambda & i & 0\\
0 & \frac{(1-\lambda)^2 + 1}{1-\lambda} & i\\
0 & i & 1-\lambda
}
\xleftrightarrow{R_3 \to R_3 - \frac{i(1-\lambda)}{(1-\lambda)^2 + 1} R_2}
\myvec{
1-\lambda & i & 0\\
0 & \frac{(1-\lambda)^2 + 1}{1-\lambda} & i\\
0 & 0 & \frac{(1-\lambda)^3 + 2(1-\lambda)}{(1-\lambda)^2 + 1}
}
\end{align}

Now determinant = product of diagonal entries:
\begin{align}
(1-\lambda)\Big(\frac{(1-\lambda)^2 + 1}{1-\lambda}\Big)\Big(\frac{(1-\lambda)^3 + 2(1-\lambda)}{(1-\lambda)^2 + 1}\Big) = 0
\end{align}

Simplifying,
\begin{align}
(1-\lambda)\big((1-\lambda)^2 + 2\big) = 0
\end{align}

Hence eigenvalues are:
\begin{align}
\lambda_1 = 1,\ 
\lambda_2 = 1 + i\sqrt{2},\ 
\lambda_3 = 1 - i\sqrt{2}
\end{align}

\textbf{Conclusion :}
For a complex symmetric matrix,
\begin{enumerate}
  \item Eigenvalues may be complex.
  \item Eigenvectors corresponding to distinct eigenvalues are not necessarily orthogonal.
\end{enumerate}

We solve \((\vec{A}-\lambda\vec{I})\vec{v}=\vec{0}\) for each eigenvalue.

\bigskip

\textbf{Eigenvalue } \(\lambda_1 = 1\).

\begin{align}
\vec{v}_1 &= \myvec{1\\0\\-1}
\end{align}

\bigskip

\textbf{Eigenvalue } \(\lambda_2 = 1 + i\sqrt{2}\).

\begin{align}
\vec{v}_2 &= \myvec{1\\[2pt]\sqrt{2}\\[2pt]1}
\end{align}

\pagebreak

\textbf{Eigenvalue } \(\lambda_3 = 1 - i\sqrt{2}\)

\begin{align}
\vec{v}_3 &= \myvec{1\\[2pt]-\sqrt{2}\\[2pt]1}
\end{align}

\bigskip

\textbf{Summary:}
\begin{align}
\lambda_1 &= 1,\quad \vec{v}_1 = \myvec{1\\0\\-1},\\
\lambda_2 &= 1 + i\sqrt{2},\quad \vec{v}_2 = \myvec{1\\\sqrt{2}\\1},\\
\lambda_3 &= 1 - i\sqrt{2},\quad \vec{v}_3 = \myvec{1\\-\,\sqrt{2}\\1}.
\end{align}

Some of the eigen values are complex but the eigen vectors are real.

\end{document}
